% Gesamttest «Kreis und Kreisteile» — Quelle des PDFs (python3 scripts/build-lp-pdf.py kreis-kreisteile).
% Musterlösung und Raster: bewertungspaket.tex. Alle Werte mit python3 nachgerechnet (scripts/lp/kreis-kreisteile/zahlen.py, 08.10.2026).
% Kein \lt/\gt — pdfLaTeX kennt sie nicht, < und > tun es.
\documentclass[11pt]{article}
\usepackage{../lp-druck}
\renewcommand{\lpname}{Kreis und Kreisteile}
\renewcommand{\lpteilgebiet}{5.2c}
\renewcommand{\lpbereich}{Grundlagenfach}
\renewcommand{\lpfuss}{Gesamttest · 25 Punkte}
\begin{document}
\lpkopf{Gesamttest}{%
  \vspace{4pt}\noindent\begin{tabularx}{\linewidth}{@{}X X@{}}
  Code (kein Name): \hrulefill & Punkte: \hrulefill\ / 25
  \end{tabularx}}

\begin{lpinfo}{Anleitung}
\textbf{Zeit:} rund 30 Minuten. \textbf{Hilfsmittel:} Taschenrechner erlaubt, für alle Aufgaben (RLP 5.2 trägt keinen Vermerk
«auch ohne Hilfsmittel»); rechne mit der $\pi$-Taste. Mach zu jeder Sachaufgabe eine \textbf{Skizze} und schreib den
\textbf{Rechenweg} auf — bewertet wird der Weg. Längen, Flächen und Winkel auf zwei Dezimalen runden, Zwischenresultate ungerundet
weiterverwenden. Figuren sind Skizzen, nicht massstäblich.
Danach: Lösung scannen oder fotografieren (ohne Namen und Standort) und mit dem \emph{Bewertungspaket} bewerten lassen.
\end{lpinfo}

\lpteil{Aufgaben G1 bis G7}{Kapitel 1 bis 4 · Taschenrechner erlaubt · 25 Punkte}

\begin{lpaufgabe}{G1}{Gerade und Kreis}{3 P}
Ein Kreis hat den Durchmesser $d = 8\,\text{cm}$. Drei Geraden $g$, $h$ und $k$ haben vom Mittelpunkt $M$ die Abstände
$a_g = 4\,\text{cm}$, $a_h = 2.5\,\text{cm}$ und $a_k = 5\,\text{cm}$.
\begin{enumerate}[label=(\alph*),leftmargin=*,itemsep=2pt]
\item Ist $g$, $h$ und $k$ je eine Passante, eine Tangente oder eine Sekante? Begründe kurz.
\item Die Gerade $h$ schneidet aus dem Kreis eine Sehne heraus. Wie lang ist diese Sehne?
\end{enumerate}
\lpplatz{5}
\end{lpaufgabe}

\begin{lpaufgabe}{G2}{Wie weit reicht der Blick?}{3 P}
Die Lampe eines Leuchtturms steht $50\,\text{m}$ über dem Meer. Ein Lichtstrahl, der das Meer gerade noch streift, berührt die
Erdkugel (Radius $6371\,\text{km}$). Wie weit ist dieser Berührpunkt von der Lampe entfernt? Zeichne eine Skizze mit dem
Erdmittelpunkt (im Querschnitt ist die Erde ein Kreis).
\lpplatz{6}
\end{lpaufgabe}

\begin{lpaufgabe}{G3}{Teich mit Weg}{4 P}
Ein kreisrunder Teich hat den Umfang $25\,\text{m}$.
\begin{enumerate}[label=(\alph*),leftmargin=*,itemsep=2pt]
\item Wie gross ist seine Fläche?
\item Rund um den Teich wird ein $1.5\,\text{m}$ breiter Weg angelegt. Wie gross ist die Fläche des Wegs?
\end{enumerate}
\lpplatz{8}
\end{lpaufgabe}

\begin{lpaufgabe}{G4}{Ein Sektor, rückwärts}{4 P}
Ein Kreissektor mit dem Radius $r = 7.5\,\text{cm}$ hat die Fläche $50\,\text{cm}^2$.
\begin{enumerate}[label=(\alph*),leftmargin=*,itemsep=2pt]
\item Wie gross ist sein Zentriwinkel $\varphi$?
\item Wie lang ist sein Bogen $b$?
\item Wie lang ist sein ganzer Rand?
\end{enumerate}
\lpplatz{9}
\end{lpaufgabe}

\begin{lpaufgabe}{G5}{Ein Segment}{4 P}
In einem Kreis ist eine Sehne $10\,\text{cm}$ lang und hat vom Mittelpunkt den Abstand $5\,\text{cm}$.
\begin{enumerate}[label=(\alph*),leftmargin=*,itemsep=2pt]
\item Berechne den Radius des Kreises.
\item Begründe, dass der Zentriwinkel über dieser Sehne $90^\circ$ beträgt.
\item Berechne die Fläche des kleineren Segments, das die Sehne abschneidet.
\end{enumerate}
\lpplatz{7}
\end{lpaufgabe}

\begin{lpaufgabe}{G6}{Eine fremde Aussage}{3 P}
Noah behauptet: «Verdoppelt man bei einem Kreissektor den Radius \emph{und} den Zentriwinkel, wird die Sektorfläche viermal so
gross.» Stimmt das? Begründe mit der Formel für die Sektorfläche. Was geschieht dabei mit der Bogenlänge?
\lpplatz{5}
\end{lpaufgabe}

\begin{lpaufgabe}{G7}{Sportplatz}{4 P}
\begin{minipage}[t]{0.52\linewidth}\vspace{0pt}
Ein Sportplatz besteht aus einem Rechteck, $100\,\text{m}$ lang und $64\,\text{m}$ breit, und zwei Halbkreisen an den beiden
schmalen Seiten (Bild).
\begin{enumerate}[label=(\alph*),leftmargin=*,itemsep=2pt]
\item Wie lang ist sein Rand (die innerste Laufbahn)?
\item Wie gross ist seine Fläche?
\end{enumerate}
\end{minipage}\hfill
\begin{minipage}[t]{0.44\linewidth}\vspace{0pt}\centering
\begin{tikzpicture}[scale=0.032]
  \fill[lpblau!8] (0,0) rectangle (100,64);
  \fill[lpblau!8] (100,0) arc (-90:90:32) -- cycle;
  \fill[lpblau!8] (0,64) arc (90:270:32) -- cycle;
  \draw[lpblau,thick] (0,0) -- (100,0) arc (-90:90:32) -- (0,64) arc (90:270:32) -- cycle;
  \draw[lpgrau,dashed] (0,0) -- (0,64); \draw[lpgrau,dashed] (100,0) -- (100,64);
  \node[below] at (50,0) {\small $100\,\text{m}$};
  \node[left] at (100,32) {\small $64\,\text{m}$};
\end{tikzpicture}
\end{minipage}
\lpplatz{6}
\end{lpaufgabe}
\end{document}
